\chapter{Cấu hình chi tiết và quy trình tái lập XBone-Net}
\label{AppendixModelDetails}

Phụ lục này bổ sung các thông tin triển khai cần thiết để tái lập phương pháp được trình bày trong Chương~\ref{Chapter3}. Nội dung tập trung vào định danh mô hình nền tảng, kích thước các biểu diễn, vị trí thích nghi hạng thấp, cấu hình dung hợp, đầu phân lớp và thứ tự huấn luyện. Các quy tắc tiền xử lý ảnh và công thức tối ưu vẫn được giữ tại Chương~\ref{Chapter3} để tránh lặp lại phần mô tả phương pháp.

\section{Mô hình nền tảng và chuẩn hóa đầu vào}

XBone-Net sử dụng trọng số công bố của BiomedCLIP trên trang web hugging face có tên là BiomedCLIP-PubMedBERT-256-vit-base-patch16-224 \cite{BiomedCLIP}. Bộ trọng số này này xác định đồng thời bộ mã hóa thị giác ViT-B/16, bộ mã hóa văn bản PubMedBERT, hai lớp chiếu về không gian 512 chiều, bộ biến đổi ảnh và bộ tách đơn vị biểu diễn tương ứng. Trọng số nền của hai bộ mã hóa được nạp trước khi chèn LoRA và được giữ cố định trong cả hai giai đoạn huấn luyện.

\begin{table}[H]
\centering
\small
\resizebox{\textwidth}{!}{%
\begin{tabular}{|L{3.4cm}|L{5.0cm}|L{6.0cm}|}
\hline
\multicolumn{1}{|C{3.4cm}|}{\textbf{Thành phần}} &
\multicolumn{1}{C{5.0cm}|}{\textbf{Cấu hình}} &
\multicolumn{1}{C{6.0cm}|}{\textbf{Vai trò}} \\ \hline
Ảnh toàn cục & Một ảnh $224\times224$ & Giữ bố cục giải phẫu chung của ảnh X-quang. \\ \hline
Ảnh cục bộ & Bốn vùng $224\times224$ được cắt từ ảnh gốc & Giữ chi tiết tại các vùng được quy trình toàn cục--cục bộ lựa chọn. \\ \hline
Mảnh ảnh & Lưới $14\times14$, mỗi mảnh $16\times16$ & Tạo 196 đơn vị biểu diễn không gian và một đơn vị đại diện cho mỗi ảnh. \\ \hline
Văn bản & Tối đa 256 đơn vị biểu diễn & Chuỗi dài hơn bị cắt, chuỗi ngắn hơn được đệm và mặt nạ chú ý loại các vị trí đệm. \\ \hline
Văn bản khuyết & no clinical information available. & Là chuỗi thay thế cố định khi trường văn bản rỗng. \\ \hline
Chiều biểu diễn & $D=512$ & Là chiều chung của ảnh, văn bản, dung hợp và tâm lớp. \\ \hline
\end{tabular}%
}
\caption{Cấu hình đầu vào và mô hình nền tảng của XBone-Net}
\label{tab:appendix-model-input-config}
\end{table}

Văn bản được đọc bằng mã hóa UTF-8, loại khoảng trắng ở hai đầu và chuyển về chữ thường trước khi tách đơn vị biểu diễn. Đối với CTCH, đầu vào gồm lý do vào viện, diễn tiến bệnh và tiền sử. Đối với BTXRD, đầu vào là văn bản tổng hợp đã được cố định trước khi huấn luyện. Bộ tách đơn vị biểu diễn và phép chuẩn hóa ảnh được lấy trực tiếp từ cùng định danh BiomedCLIP nhằm tránh sai khác giữa tiền huấn luyện và thích nghi trên dữ liệu đích.

\section{Kích thước biểu diễn qua các thành phần}

Gọi $B$ là kích thước lô và $K$ là số lớp. Cấu hình chính sử dụng $K=22$ trên CTCH và $K=10$ trên BTXRD. Bảng~\ref{tab:appendix-model-tensor-shapes} trình bày kích thước đầu ra tại các bước chính. Các kích thước này không phụ thuộc vào độ phân giải ban đầu của ảnh vì ảnh toàn cục và các vùng cục bộ đều được đưa về $224\times224$ trước khi mã hóa.

\begin{table}[H]
\centering
\small
\resizebox{\textwidth}{!}{%
\begin{tabular}{|L{4.4cm}|C{4.2cm}|L{6.0cm}|}
\hline
\multicolumn{1}{|C{4.4cm}|}{\textbf{Vị trí}} &
\textbf{Kích thước} &
\multicolumn{1}{C{6.0cm}|}{\textbf{Diễn giải}} \\ \hline
Chuỗi của một ảnh sau ViT & $B\times197\times512$ & Một đơn vị đại diện và 196 đơn vị mảnh ảnh. \\ \hline
Một vùng ảnh cục bộ sau gộp & $B\times5\times512$ & Một đơn vị đại diện và lưới gộp $2\times2$. \\ \hline
Bốn vùng ảnh cục bộ & $B\times20\times512$ & Ghép năm đơn vị của mỗi vùng theo thứ tự xác định. \\ \hline
Chuỗi thị giác dùng cho dung hợp & $B\times21\times512$ & Một biểu diễn ảnh toàn cục và 20 đơn vị ảnh cục bộ đã bổ sung tọa độ. \\ \hline
Chuỗi văn bản & $B\times256\times512$ & Một đơn vị đại diện, các đơn vị nội dung và các vị trí đệm. \\ \hline
Biểu diễn ảnh nhận ngữ cảnh văn bản & $B\times512$ & Kết quả truy vấn văn bản bằng biểu diễn ảnh toàn cục. \\ \hline
Biểu diễn văn bản nhận ngữ cảnh ảnh & $B\times512$ & Kết quả truy vấn các đơn vị ảnh cục bộ bằng biểu diễn văn bản toàn cục. \\ \hline
Biểu diễn dung hợp & $B\times512$ & Đầu vào của bộ phân lớp theo tâm lớp và các điểm OOD hậu xử lý. \\ \hline
Logit phân lớp & $B\times K$ & Điểm của từng lớp trước Softmax. \\ \hline
\end{tabular}%
}
\caption{Kích thước các biểu diễn tại những bước chính của XBone-Net}
\label{tab:appendix-model-tensor-shapes}
\end{table}

Trong pha 1, bốn vùng cục bộ được tổng hợp bằng trọng số chú ý có điều kiện theo biểu diễn ảnh toàn cục. Bản tóm tắt cục bộ được nhân với hệ số $0{,}25$ trước khi cộng với biểu diễn toàn cục và chuẩn hóa. Trong pha 2, cấu hình chính không nén 20 đơn vị ảnh cục bộ bằng truy vấn học được mà truyền trực tiếp toàn bộ chuỗi này vào mô-đun dung hợp.

\section{Cấu hình thích nghi hạng thấp}

LoRA \cite{Hu2022LoRA} được chèn vào cả bộ mã hóa ảnh và bộ mã hóa văn bản. Cấu hình chung sử dụng hạng $r=16$, hệ số tỷ lệ $\alpha=32$ và tỷ lệ bỏ học $0{,}1$. Các ma trận LoRA được cập nhật trong cả hai pha, còn các trọng số nền tương ứng không nhận gradient. Bảng~\ref{tab:appendix-model-lora-targets} liệt kê các phép chiếu tuyến tính được thích nghi.

\begin{table}[htbp]
\centering
\resizebox{\textwidth}{!}{%
\begin{tabular}{|l|l|l|}
\hline
\multicolumn{1}{|c|}{\textbf{Bộ mã hóa}} &
  \multicolumn{1}{c|}{\textbf{Mô-đun đích}} &
  \multicolumn{1}{c|}{\textbf{Ý nghĩa}} \\ \hline
Thị giác ViT-B/16 &
  \begin{tabular}[c]{@{}l@{}}attn.qkv, attn.proj,\\ mlp.fc1, mlp.fc2\end{tabular} &
  \begin{tabular}[c]{@{}l@{}}Thích nghi các phép chiếu \\ truy vấn--khóa--giá trị, đầu \\ ra chú ý và hai lớp của mạng \\ truyền thẳng trong mỗi khối \\ thị giác.\end{tabular} \\ \hline
Văn bản PubMedBERT &
  \begin{tabular}[c]{@{}l@{}}query, key, value,\\ attention.output.dense,\\ intermediate.dense,\\ output.dense\end{tabular} &
  \begin{tabular}[c]{@{}l@{}}Thích nghi các phép chiếu tự \\ chú ý, đầu ra chú ý và hai phép \\ chiếu của mạng truyền thẳng \\ trong mỗi khối văn bản.\end{tabular} \\ \hline
\end{tabular}%
}
\caption{Các mô-đun tuyến tính được thích nghi bằng LoRA}
\label{tab:appendix-model-lora-targets}
\end{table}

Các tham số học được ngoài LoRA gồm phép chiếu tọa độ của vùng ảnh, hệ số tọa độ khởi tạo bằng $0{,}1$, hàm chấm điểm dùng để tổng hợp ảnh cục bộ, mô-đun chú ý chéo, mạng dung hợp và độ lệch theo lớp của đầu phân loại. Các tâm lớp là thống kê của biểu diễn tập huấn luyện nên không được tối ưu trực tiếp bằng lan truyền ngược.

\section{Dung hợp và đầu phân lớp}

Mô-đun dung hợp sử dụng chú ý chéo hai chiều với tám đầu chú ý, chiều ẩn 512 và tỷ lệ bỏ học $0{,}1$. Ở chiều thứ nhất, biểu diễn ảnh toàn cục truy vấn các đơn vị văn bản không đệm. Ở chiều còn lại, biểu diễn văn bản toàn cục truy vấn 20 đơn vị ảnh cục bộ. Hai kết quả được cộng thặng dư với biểu diễn toàn cục tương ứng, chuẩn hóa lớp, nối với nhau và đưa qua mạng hai lớp để tạo biểu diễn dung hợp 512 chiều.

\begin{table}[H]
\centering
\small
\resizebox{\textwidth}{!}{%
\begin{tabular}{|L{4.0cm}|L{4.0cm}|L{6.5cm}|}
\hline
\multicolumn{1}{|C{4.0cm}|}{\textbf{Thành phần}} &
\multicolumn{1}{C{4.0cm}|}{\textbf{Giá trị}} &
\multicolumn{1}{C{6.5cm}|}{\textbf{Ghi chú}} \\ \hline
Hướng chú ý & Hai chiều & Ảnh truy vấn văn bản và văn bản truy vấn ảnh dùng hai mô-đun độc lập. \\ \hline
Số đầu chú ý & 8 & Mỗi đầu có chiều $512/8=64$. \\ \hline
Tỷ lệ bỏ học & $0{,}1$ & Áp dụng trong chú ý, nhánh thặng dư và mạng dung hợp. \\ \hline
Mạng dung hợp & $1024\rightarrow512\rightarrow512$ & Sử dụng GELU giữa hai lớp tuyến tính và chuẩn hóa đầu ra. \\ \hline
Hệ số logit của đầu tâm lớp & $15{,}0$ & Nhân với độ tương đồng cosine trước khi cộng độ lệch lớp. \\ \hline
Độ lệch theo lớp & Có, khởi tạo bằng 0 & Được tối ưu trong pha 2 để hiệu chỉnh ranh giới quyết định theo lớp. \\ \hline
Cập nhật tâm lớp & Mỗi vòng huấn luyện & Tính lại từ toàn bộ tập huấn luyện trước khi đánh giá trên tập xác thực. \\ \hline
\end{tabular}%
}
\caption{Cấu hình mô-đun dung hợp và đầu phân lớp}
\label{tab:appendix-model-fusion-classifier}
\end{table}

\section{Cấu hình huấn luyện và lựa chọn trạng thái mô hình}

Hai pha sử dụng AdamW, lịch tốc độ học cosine, tỷ lệ khởi động $0{,}1$, hệ số suy giảm trọng số $10^{-4}$ và kích thước lô 16. Huấn luyện sử dụng bfloat16 trên ba hạt giống 42, 123 và 456. Pha 1 lựa chọn trạng thái có mất mát xác thực thấp nhất, còn pha 2 lựa chọn trạng thái có Macro-F1 xác thực cao nhất. Tập kiểm thử không tham gia dừng sớm, lựa chọn trạng thái hoặc cập nhật tâm lớp.

\begin{table}[H]
\centering
\small
\resizebox{\textwidth}{!}{%
\begin{tabular}{|L{3.5cm}|C{3.1cm}|C{3.1cm}|L{5.2cm}|}
\hline
\multicolumn{1}{|C{3.5cm}|}{\textbf{Thuộc tính}} &
\textbf{Pha 1} &
\textbf{Pha 2} &
\multicolumn{1}{C{5.2cm}|}{\textbf{Diễn giải}} \\ \hline
Mục tiêu & Khớp ngữ nghĩa & Phân lớp theo tâm lớp & Pha 1 căn chỉnh ảnh--văn bản, pha 2 tối ưu biểu diễn dung hợp cho lớp đích. \\ \hline
Tốc độ học & $5\times10^{-5}$ & $10^{-4}$ & Áp dụng cho các tham số được phép cập nhật trong từng pha. \\ \hline
Số vòng tối đa & 30 & 50 & Đánh giá và lưu trạng thái sau mỗi vòng. \\ \hline
Dừng sớm & 3 vòng & 5 vòng & Dừng khi độ đo lựa chọn không cải thiện. \\ \hline
Độ đo lựa chọn & Mất mát xác thực & Macro-F1 xác thực & Chiều tối ưu lần lượt là giảm và tăng. \\ \hline
Đích mềm cùng lớp & $\rho=0{,}50$ & Không áp dụng & Giảm việc xem cặp khác mẫu nhưng cùng lớp là âm hoàn toàn. \\ \hline
Làm trơn nhãn & Không áp dụng & $\varepsilon=0{,}05$ & Giới hạn việc phân phối mục tiêu tập trung tuyệt đối vào một lớp. \\ \hline
Trọng số lớp & Không áp dụng & Số lượng hiệu dụng & Giới hạn trọng số lớn nhất bằng $5{,}0$. \\ \hline
\end{tabular}%
}
\caption{Cấu hình và tiêu chí lựa chọn trạng thái trong hai pha huấn luyện}
\label{tab:appendix-model-training-phases}
\end{table}

Hai bộ dữ liệu chỉ khác nhau ở số lớp và hệ số số lượng hiệu dụng. CTCH sử dụng $K=22$ và $\beta=0{,}999$, trong khi BTXRD sử dụng $K=10$ và $\beta=0{,}99$. Các cấu hình còn lại trong Bảng~\ref{tab:appendix-model-training-phases} được giữ giống nhau. Trong thiết lập ít mẫu, chỉ tập huấn luyện được lấy tối đa 1, 10 hoặc 20 mẫu trên mỗi lớp theo từng hạt giống. Tập xác thực và tập kiểm thử được giữ nguyên.